\startmode[single]

% \enabletrackers[publications, publications.crossref, publications.details, publications.cite, publications.strings]
%
% \usebtxdataset[main][bib/default.bib,bib/physik.bib,bib/mathe.bib,bib/seminarfach.bib]
% \usebtxdefinitions[aps]
% \setupbtx[dataset=main]
% \definebtxrendering[bibrendering][aps][dataset=main]

\setvariables[meta]
  [title={PH10_AB_Nuklidkarte},
   subtitle={},
   date={2025-03-14},
   author={Niklas von Hirschfeld},
  ]

\setupinteraction[title={\getvariable{meta}{title}}]

\starttext

\startfrontmatter
\component c_titlepage
\component c_toc
\stopfrontmatter

\setuppagenumber[number=1]

\startbodymatter

\startchapter[title=Nuklidkarte]

\stopmode

\section{Kennenlernen der Nuklidkarte}

Zusätzlich angegeben:

\startitemize[packed,m]
\item Die Energie in MeV: $\alpha$
\item Die Ordnungszahl: $Z$
\item Die Anzahl der Neutronen: $N$
\stopitemize

\blank
\hrule
\blank

Man verliert zwei Neutronen und zwei Protonen. Mann muss also zwei nach links
(Neutrone) und zwei nach unten (Protonen). Damit landet man auf {\bf Rn222}.
Ra-226 hat $\alpha -$ Starhlung, welche aus zwei Protonen und zwei Neutronen
besteht.

\blank
\hrule
\blank

Bi-210. $\beta -$Strahlung. Wir verlieren ein Neutron und bekommen ein Proton
(ein nach links und ein nach oben). Also ist das Nuklid: Po210

\startframedtipp
Nuklidkarte benuzten:

\startplacetable[location={here},title={}]
\bTABLE
\startCSV
$\alpha$, 2 links, 2 runter
$\beta^-$, 1 links, 1 hoch
$\gamma$, 
\stopCSV
\eTABLE
\stopplacetable
\stopframedtipp


\section{Analyse des radioaktiven Mülls}

\startframedtipp
Castor behälter 
\stopframedtipp



\startplacetable[location={here},title={Nuklide für den CASTOR}]
\bTABLE
\startCSV
Nuklid, Strahlungsart, Halbwertszeit

\stopCSV
\eTABLE
\stopplacetable


\startmode[single]

\stopchapter

\stopbodymatter

\startbackmatter
% \component c_definitions
\component c_bibliography
\stopbackmatter
\stoptext

\stopmode
